\subsection{Προελεύση και χοροχρονική κατανομή των εξωτερικών ρύπων}
Υπάρχουν πολλές πηγές που εκπέμπουν ρύπους στην ατμόσφαιρα, τόσο φυσικές όσο και ανθρωπογενείς, με τις ανθρωπογενείς να είναι οι κυρίαρχες από την αρχή της βιομηχανοποίησης. Η καύση είναι μεγαλύτερος συντελεστής ρύπανσης. 

Εξωτερικές πηγές καύσης αποτελούν την χρήση ορυκτών καυσίμων για  τα διάφορα είδη μεταφοράς, την βιομηχανία και την παραγωγή ηλεκτρικής ενέργειας, την καύση βιομάζας, οπού συμπεριλαμβάνει την ελεγμένη ή όχι φωτιά σε δάση και σαβάνα, η καύση γεωργικών αποβλήτων όπως και η καύση σκουπιδιών στις πόλεις. Πηγές εκπομπής ρύπων που δεν προέρχονται από άμεση κάψη στον εξωτερικό χώρο είναι η ανασοίκωση της σκόνης,κατασκευαστικές διαδικάσιες, όπως και μεταφορά της ρύπανσης από ένα μέρο σε ένα άλλο, κάτι που απασχολεί κατά κύριο λόγο την αστική ρύπανση.

Οι ρύποι  κατηγοροποιούνται σε προτωγενής και δευτερογενής με βάση εάν εκπέμπονται κατευθείαν στον αέρα από την πηγή προέλευσης του, ή δημιουργούνται από φυσικοχημικές διαδικασίες που υπόκοινται οι πρωτογενής ρύποι ενώ βρίσκονται είδη στον αέρα , Η εκπομπή μικροσωματιδίων που προέρχονται από την καύση είναι προτωγενής ρύπος, ενώ όταν το οξίδια του αζώτου και μη πτητικές ουσίες δημιουργούν όζον με χημική αντίδραση στην υψηλή θερμοκρασία και στην φωτεινότητα, είναι δευτερογενής ρύπος \cite{EPA-GROUND-OZONE}.


Η χωροχρονική συγκέντρωση των ρύπων εξαρτάται από που βρίσκονται οι πηγές και το πρόγραμμα χρονικής λειτουργίας τους (καθημερινή ή εποχική), τα χαρακτηριστικά των ρύπων και τις δυναμικές συμπεριφορές τους (διασπορά,εναπόθεση, αλλιλεπίδραση με άλλους ρύπους) και τις μετερολογικές συνθήκες.

Στα αστικά περιβάλλοντα, κάποιοι ρύποι διασκορποίζονται πιο ομοιογενώς από άλλους. Για παράδειγμα η οικογένεια των μικροσωματιδίων που έχουν διάμετρο 2,5 μικρόμετρα και κάτω (συνήθως αναφερόμενο εώς PM\textsubscript{2.5} , 
έχουν πολύ μικρότερη ΄χωρική διακύμανση των συγκεντρώσεων τους σε μία δεδομένη έκταση από ότι έχουν τα υπέρελεπτα σωματίδια (ultra fine particles ή UFP), όπου είναι η οικογένεια των μικροσωματιδίων με διάμετρο κάτω από 0.1 μικρόμετρα ή  αέριοι ρύποι που διασκορπούνται κατευθείαν από μια ρυπογόνα πηγή.



Ο λόγος που είναι σημαντική η χωρική διακύμανση των ρύπων σε μία περίοχη είναι κατά πόσο αποτελεσμάτικη είναι μία καταμέτρηση των συγκεντρώσεων σε κάποια συγκεκριμένη σταθερή τοποθεσία να αποδώσει την ατμοσφαίρικη ρύπανση που υπάρχει σε διαφορετικά σημεία της περιόχης αυτής. 

Η χρονική διακύμανση είναι επίσης πολύ σημαντικό στοιχεία γιατί μας δίνει βαθύτερες κατανοήσεις για τα είδικα χρονικά μότιβα που μπορεί να έχουν οι ρύποι και μας δίνει την δυνατότητα προβλέψης της συμπεριφοράς τους (πχ πως αλλάζει η ρύπανση που προκαλούν τα αυτοκίνητα τις καθημερινές σε σχέση με τα σαββατοκύριακα). Βεβαίως, και η κατάσταση του καιρού είναι ένας πολύ σημαντικός παράγοντας των χρονικών μεταβάσεων της συγκεντρώσεων των ρύπων, πολύ μεγαλύτερος ακόμα και από την ίδια την χρόνικη μετάβαση που παρουσιάζει η ίδια η πηγή στην εκπομπή των ρύπων.


\subsection{Προελεύση και χοροχρονική κατανομή των εσωτερικών ρύπων}

Στα εσωτερικά περιβάλλοντα, η καύση προέρχονται συνήθως είτε από την  χρήση βλαβερών καυσίμων για θέρμανση και μαγείρευμα χωρίς κατάλληλο εξαερισμό, όπως είναι το κάρβουνο, το ξύλο και η κοπριά, είτε από μικρότερα είδη καύσης όπως κεριά,θυμιάματα και λάμπες κηροζίνης. Το κάπνισμα αποτελεί επίσης ενά πολύ σημαντικό κομμάτι ρύπων που δεν πρέπει να ξεχνάμε.

Στον εσωτερικό αέρα, οι πηγές και διαδικασίες που δεν προέρχονται από καύση είναι πολύ πιο κυρίαρχες, συγκεκριμένα πηγές που παραγούν πτητικές και ημηπτιτικές οργανικές πηγές ρύπων (Volatile και semi-volatile organic compounds). Αυτές περιλαμβάνουν πτητικές οργανικές ενώσεις που εξάγωνται κατά τον πρώτο καιρό μετά από κατασκευή ή ανακαίνωση ενος κτηρίου ή χώρου, χρήση καταναλοτικών προίοντων (παρασιτοκτώνα και καθαριστικά προιόντα), όπως και από την λειτουργία ηλεκτρονικών συσκευών όπως εκτυπωτές. Ακόμα,η κίνηση και δραστηριότητα των ανθρώπων δημιουργία την επαναιώριση της σκόνης είναι μία ακόμα σημαντική πηγή ρύπων, ειδικά στα σχολεία.

Παρόλα αυτά, οι ρύποι του εσωτερικού αέρα δεν προέρχονται μόνο από τον εσωτερικές πηγές αλλά και από ρύπους που έρχονται από τον εξωτερικό αέρα, με την Οι διόδος τους μα είναι Ο εξαερισμός και η δυισδισή του περιβλήματος του κτηρίου. Κάποιοι από αυτούς τους ρύπους αποσυντίθενται γρήγορα στον εσωτερικό χώρο, όπως το όζον, ενώ αλλά είναι αδρανείς σε αλλαγές, όπως το μονοξείδιο του αζώτου.

Η συγκέντρωση τους εξαρτάται πολύ από την απόσταση του κτηρίου από τις εξωτερικές πηγές και από τα χρονικά μοτίβα που κυριαρχούν, τον ρυθμό εξαερισμού, τον βαθμό διείσδισης του κτηρίου και τον ρυθμό αποσύνθεσης τον ρύπων αυτών.
Άρα ακόμα και χωρίς εσωτερικές πηγές ρύπων, μπορεί να υπάρχει ρύπανση, η οποία θα προέρχεται από τον εξωτερικό χώρο. Η κατηγοροιοπήση εσωτερικών και εξωτερικών ρύπων υπάρχουν και στον εσωτερικό χώρο, αλλά είναι πιο δύσκολο να μελετηθούν \cite{SARAGA2024117821}.


\section{Σημασία μελέτης ποιότητας εσωτερικού αέρα}

Οι άνθρωποι εκτίθονται σε επαφή με ρύπους σε κάθε μικροπεριβάλλον που παιρνάνε κάποιο χρόνο και η έκθεση αυτή τους βάζει σε ένα βαθμό ρίσκου \cite{WHO(2021)}. Ένα μικροπεριβάλλον ορίζεται σαν ένας τρισδιάστατος χώρος ο οποίος έχει ομοιόμορφή συγκέντρωση ρύπων σε μία δεδομένη χρονική στιγμή. Η εκθεσή που έχει ένας άνθρωπος σε ένα ρύπο είναι προιόν όχι μόνο των συγκεντρόσεων του ρύπου σε ένα συγκεκριμένο μικρο περιβάλλον αλλά και ο χρόνος που θα περάσει στο μικροπεριβάλλον αυτό.

Οπότε ενώ σε εσωτερικούς χώρους όπου δεν υπάρχει καύση βλαβερών καυσίμων, η ρύπανση να είναι σχετικά μικρή σε σχέση με χώρους με μεγάλη ρύπανση (πχ δρόμοι με μεγάλη κίνηση), ο χρόνος που περνάμε στους εσωτερικούς χώρους αποτελεί σημαντικό συντελεστή της ποιότητας υγείας. Και οι άνθρωποι σήμερα παιρνάνε πολύ χρόνο μέσα.

Σύμφωνα με μελετές, %βρες τις μελετές από διάφορα μέρη
οι άνθρωποι παιρνάνε πάνω από το 80 με 90 της εκατό της ημέρας στους είτε μέσα σε κάποιο κτήριο είτε σε κάποιο μεταφορικό μέσο .

\section{Όροι που χρησιμοποιούνται}

\subsection{Πρωτογενής και δευτερογενής ρύποι}

\subsection{μετά από εδώ δεν έχω αρχίσει το γράψιμο}

Σύμφωνα με την United States Environmental Protection Agency (EPA) \cite{EPA(1991)} 

%to kato tha paei allou
Και ενώ το κάψη στεραίων καυσίμων χωρίς τον κατάλληλο εξαερισμό έχει υποχωρήσει πολύ σημαντικά στις οικονομικά ανεπτυγμένες χώρες, σε μεγάλο μέρος του πλανήτη συνεχίζει η καύση υλών 

